\section{Sběrnice}

Kromě propojovacího prvku, sběrnice slouží také jako správní centrum. Obsahuje frontu událostí, jako
např. požadavek o načtení programu z SD karty nebo resetování, které zpracuje a odpovídajícím
způsobem aktualizuje stav. Je velice riskantní provádět tento požadavek hned v okamžik přijetí
požadavku, jelikož hrozí to poškozením stavu emulátoru. Proto byla navržená fronta události, která
pomáhá odložit požadavky a zpracovat je v jeden, správný okamžik.

Když CPU přistupuje ke sběrnici, využívá funkce \texttt{bus\_read()} a \texttt{bus\_write()}, které
realizují mapování systému NES.

\begin{listing}[htbp]
\begin{minted}{c}
void bus_write(Bus *self, uint16_t addr, uint8_t data) {
    /* Optimalizace: místo žebříku ifů je efektivnější rozdělit
     * adresový prostor na 4KB stránky */
    switch (addr >> 12) {
    case 0x0:
    case 0x1:
        self->memory[addr & 0x07ff] = data;
        break;
    case 0x2:
    case 0x3:
        /* Catch-up mechanismus: při přístupu k registrům, PPU dohání stav CPU. */
        ppu_run_until(&self->ppu, self->cpu.cycles * 3);
        ppu_write_reg(&self->ppu, 0x2000 + (addr & 7), data);
        break;
    case 0x4:
        if (addr == 0x4014) {
            handle_oam_dma(self, data);
        }
        break;
        /* Ostatní mapování byly zkráceny... */
    }
}
\end{minted}
\caption{Komunikační rozhraní sběrnice (příklad funkce pro zápis do adresního prostoru).}
\end{listing}

Naivní implementaci funkcí \texttt{bus\_write()} a \texttt{bus\_read()} mohl by být žebřík ifů,
které postupně porovnávají parametr \texttt{addr} s adresovými rozsahy. Tato implementace však může
být nepřátelská vůči branch predictoru (který je obsazen i v RP2350) a také nutí procesor vykonávat
posloupnost porovnání. Není garantováno, že kompilátor může být schopen optimalizovat.

Lepší návrh je využit toho fakta, že paměť NESu je fakticky rozdělena do stránek velkých 4096 bytů
(0x1000). Například PPU je mapován do rozsahu 0x2000--0x3fff, což když vydělíme 4096, dostaneme
rozsah 2--3. Pro rychlejší dělení je možné použit logický posun \texttt{addr >> 12}.

% [vlozit kod zpracovavani eventu]

% - struktura Bus
% - každý komponent má odpovídající X_read a X_write funkce.
% - pamet je rozdelena na 4kb stranky [optimalizace]
% - použití bitových operací pro optimalizaci
% - catch-up mechanismus: pristupovany komponent dohani stav CPU
% - realizace udalosti a fronty udalosti
